\documentclass[10pt,aspectratio=169]{beamer}
\input{../../_en_preambulo_beamer}
% Comandos y macros estadísticas compartidas para las presentaciones Beamer en inglés (Metropolis)

% Conjuntos numéricos básicos
\newcommand{\R}{\mathbb{R}}
\newcommand{\N}{\mathbb{N}}
\newcommand{\Q}{\mathbb{Q}}
\newcommand{\Z}{\mathbb{Z}}
\newcommand{\set}[1]{\left\{ #1 \right\}}
\newcommand{\abs}[1]{\left|#1\right|}
\newcommand{\norm}[1]{\left\| #1 \right\|}

% Comandos estadísticos y probabilísticos del curso
\newcommand{\Var}{\operatorname{Var}}
\newcommand{\cov}{\operatorname{Cov}}
\newcommand{\corr}{\rho}
\newcommand{\comb}[2]{\begin{pmatrix} #1 \\ #2 \end{pmatrix}}
\newcommand{\s}{\sigma}
\newcommand{\particion}{\mathfrak{P}}
\newcommand{\card}[1]{\#\left( #1 \right)}
\newcommand{\seq}[3]{\set{#1}_{#2}^{#3}}
\newcommand{\E}{\mathbb{E}}
\newcommand{\Prob}{\mathbb{P}}

% Entornos pedagógicos y cajas en inglés académico natural (soportando tanto nombres en inglés como en español)
\newenvironment{discussion}{%
  \begin{alertblock}{Discussion Question}%
}{%
  \end{alertblock}%
}
\newenvironment{discusion}{%
  \begin{alertblock}{Discussion Question}%
}{%
  \end{alertblock}%
}

\newenvironment{reflection}{%
  \begin{alertblock}{Classroom Reflection}%
}{%
  \end{alertblock}%
}
\newenvironment{reflexion}{%
  \begin{alertblock}{Classroom Reflection}%
}{%
  \end{alertblock}%
}

\newenvironment{paradox}{%
  \begin{alertblock}{Exploratory Paradox}%
}{%
  \end{alertblock}%
}
\newenvironment{paradoja}{%
  \begin{alertblock}{Exploratory Paradox}%
}{%
  \end{alertblock}%
}

\newenvironment{motivation}{%
  \begin{exampleblock}{Motivation: Data Science in Practice}%
}{%
  \end{exampleblock}%
}
\newenvironment{motivacion}{%
  \begin{exampleblock}{Motivation: Data Science in Practice}%
}{%
  \end{exampleblock}%
}

\newenvironment{quickchallenge}{%
  \begin{block}{Quick Team Challenge}%
}{%
  \end{block}%
}
\newenvironment{retoclase}{%
  \begin{block}{Quick Team Challenge}%
}{%
  \end{block}%
}


\title[Standard Errors]{Section 6.4: Standard Errors}
\subtitle{Estimator precision and margin of error}
\author[J. Castillo Colmenares]{Juliho Castillo Colmenares}
\institute{Tecnológico de Monterrey}
\date{}

\begin{document}
\begin{frame}[plain]\vspace{-0.8cm}\titlepage\end{frame}

\begin{frame}{Roadmap --- Chapter 6}
  \footnotesize
  \begin{itemize}\itemsep=0.08cm
    \item \textbf{06.01--06.02} Point and interval estimation
    \item \textbf{06.04} Standard errors \emph{(today)}
    \item \textbf{06.05} Proportions and differences
    \item \textbf{06.07} Sample size
  \end{itemize}
  \begin{block}{Learning objective}
    Calculate theoretical standard errors and their sample estimates for means,
    differences, and proportions.
  \end{block}
\end{frame}

\begin{frame}{Definition and purpose}
  \footnotesize
  For any point estimator $\hat\theta$,
  \[
    SE(\hat\theta)=\sqrt{\operatorname{Var}(\hat\theta)}.
  \]
  \begin{alertblock}{Interpretation}
    Standard error is the standard deviation of the sampling distribution and
    the basic building block of every margin of error.
  \end{alertblock}
\end{frame}

\begin{frame}{One-population mean}
  \footnotesize
  For $\hat\theta=\bar X$,
  \[
    SE(\bar X)=\frac{\sigma}{\sqrt n},
    \qquad \widehat{SE}(\bar X)=\frac{s}{\sqrt n}.
  \]
  The second expression replaces the unknown population deviation with $s$.
\end{frame}

\begin{frame}{Independent means difference}
  \footnotesize
  For $\hat\theta=\bar X_1-\bar X_2$,
  \[
    SE=\sqrt{\frac{\sigma_1^2}{n_1}+\frac{\sigma_2^2}{n_2}},
    \quad
    \widehat{SE}=\sqrt{\frac{S_1^2}{n_1}+\frac{S_2^2}{n_2}}.
  \]
  Under equal variances, the pooled form $S_p\sqrt{1/n_1+1/n_2}$ is available.
\end{frame}

\begin{frame}{Paired means difference}
  \footnotesize
  For paired differences $D_i$ and $\hat\theta=\bar D$,
  \[
    SE(\bar D)=\frac{\sigma_D}{\sqrt n},
    \qquad \widehat{SE}(\bar D)=\frac{S_D}{\sqrt n}.
  \]
  The paired design becomes one sample of differences.
\end{frame}

\begin{frame}{Proportion and proportions difference}
  \footnotesize
  \[
    \widehat{SE}(\hat p)=\sqrt{\frac{\hat p(1-\hat p)}{n}},
  \]
  \[
    \widehat{SE}(\hat p_1-\hat p_2)=
    \sqrt{\frac{\hat p_1(1-\hat p_1)}{n_1}+
    \frac{\hat p_2(1-\hat p_2)}{n_2}}.
  \]
\end{frame}

\begin{frame}{Compact reference table}
  \footnotesize
  \begin{tabular}{lll}
    \hline
    Parameter & Estimator & Standard error \\
    \hline
    $\mu$ & $\bar X$ & $\sigma/\sqrt n$ \\
    $\mu_1-\mu_2$ & $\bar X_1-\bar X_2$ & $\sqrt{\sigma_1^2/n_1+\sigma_2^2/n_2}$ \\
    $\mu_D$ & $\bar D$ & $\sigma_D/\sqrt n$ \\
    $p$ & $\hat p$ & $\sqrt{p(1-p)/n}$ \\
    \hline
  \end{tabular}
\end{frame}

\begin{frame}{Computational bridge 1/4: sample substitution}
  \footnotesize
  When $\sigma$, $\sigma_D$, or $p$ are unknown, use $s$, $s_D$, or $\hat p$.
  \begin{block}{Principle}
    The theoretical form identifies variability; the estimated form makes it
    computable from data.
  \end{block}
\end{frame}

\begin{frame}{Computational bridge 2/4: effect of $n$}
  \footnotesize
  Means and proportions include $1/\sqrt n$.
  \begin{itemize}\itemsep=0.12cm
    \item quadrupling $n$ halves standard error;
    \item larger samples improve precision with diminishing returns;
    \item non-sampling error is not part of this formula.
  \end{itemize}
\end{frame}

\begin{frame}{Computational bridge 3/4: margin of error}
  \footnotesize
  \[
    E=(\text{critical value})\times SE.
  \]
  Standard error supplies variability; the critical value supplies confidence
  and the reference distribution.
\end{frame}

\begin{frame}{Computational bridge 4/4: formula choice}
  \footnotesize
  \begin{tabular}{ll}
    \hline
    Design & Basic estimator \\
    \hline
    One mean & $\bar X$ \\
    Two independent groups & $\bar X_1-\bar X_2$ \\
    Paired data & $\bar D$ \\
    Binary response & $\hat p$ or $\hat p_1-\hat p_2$ \\
    \hline
  \end{tabular}
\end{frame}

\begin{frame}{Derived application 1/4 --- one mean}
  \footnotesize
  The reading assigns $\bar X$ to $\mu$.
  \begin{block}{Statement}
    Obtain theoretical standard error and its estimate when $\sigma$ is unknown.
  \end{block}
\end{frame}

\begin{frame}{Derived application 1/4 --- solution}
  \footnotesize
  \[
    SE=\frac{\sigma}{\sqrt n},\qquad \widehat{SE}=\frac{s}{\sqrt n}.
  \]
  The second value is used in a sample-based interval or test.
\end{frame}

\begin{frame}{Derived application 2/4 --- two means}
  \footnotesize
  For two independent samples, uncertainty in the difference combines both
  variances.
  \begin{block}{Statement}
    Identify both sample sizes and both standard deviations.
  \end{block}
\end{frame}

\begin{frame}{Derived application 2/4 --- solution}
  \footnotesize
  \[
    \widehat{SE}=\sqrt{\frac{S_1^2}{n_1}+\frac{S_2^2}{n_2}}.
  \]
  Under equal variances, the reading permits the pooled form with $S_p$.
\end{frame}

\begin{frame}{Derived application 3/4 --- paired data}
  \footnotesize
  Each pair produces one difference variable $D$.
  \begin{block}{Statement}
    Apply the one-mean formula to the observed differences.
  \end{block}
\end{frame}

\begin{frame}{Derived application 3/4 --- solution}
  \footnotesize
  \[
    \widehat{SE}(\bar D)=\frac{S_D}{\sqrt n}.
  \]
  Dependence within each pair is absorbed into $D_i$ before precision is
  calculated.
\end{frame}

\begin{frame}{Derived application 4/4 --- proportions}
  \footnotesize
  For a Bernoulli response, the reading uses $\hat p=X/n$.
  \begin{block}{Statement}
    Estimate variability for one proportion and for the difference of two rates.
  \end{block}
\end{frame}

\begin{frame}{Derived application 4/4 --- solution}
  \footnotesize
  \[
    \widehat{SE}(\hat p)=\sqrt{\frac{\hat p(1-\hat p)}{n}}.
  \]
  For two groups, add the two estimated variances inside the square root.
\end{frame}

\begin{frame}{Synthesis of the section}
  \footnotesize
  \begin{itemize}\itemsep=0.12cm
    \item Standard error is an estimator's standard deviation.
    \item Its form depends on design and parameter.
    \item Sample substitution makes it usable with data.
    \item The margin adds a critical value.
  \end{itemize}
\end{frame}

\begin{frame}{Curricular transition}
  \footnotesize
  \begin{block}{Established}
    Precision starts by identifying the estimator and its sampling distribution.
  \end{block}
  \begin{alertblock}{Next section}
    \textbf{06.05 Proportions and Differences}: Wald, Wilson, and comparisons of
    two rates.
  \end{alertblock}
\end{frame}
\end{document}
